\subsection{Common Node Elimination}\label{sec:cne}
Common Node Elimination (CNE) permits the removal of redundant computations by detecting congruent
nodes. These nodes always produce the same results, enabling the redirection of their result edges
to a single node. This renders the other nodes dead, permitting Dead Node Elimination to remove
them. CNE is similar
to common subexpression elimination and value numbering~\cite{Alpern88} in that it detects
equivalent computations, but as the RVSDG represents all computations uniformly as nodes, it can
be extended to conditionals~\cite{Rugina00}, loops, and functions.

We consider two simple nodes $n_1$ and $n_2$ congruent, or $n_1 \cong n_2$, if they represent the
same computation, have the same number of inputs, \ie, $\card{I_{n_1}} = \card{I_{n_2}}$, and the
inputs $\iput{k}{n_1}$ and $\iput{k}{n_2}$ are congruent, or $\iput{k}{n_1} \cong \iput{k}{n_2}$,
for all $k = [0..\card{I_{n_1}}]$. Two inputs are congruent if their respective origins
$\origin{k}{n_1}$ and $\origin{k}{n_2}$ are congruent, \ie, $\origin{k}{n_1} \cong
\origin{k}{n_2}$. By definition, the origins of inputs are either outputs of simple or structural
nodes, or arguments of regions. Origins from simple nodes are only equivalent when their respective
producers are computationally equivalent, whereas for the other cases, it must be guaranteed that
they always receive the same value.

The implementation of CNE consists of two phases: mark and divert. The mark phase identifies
congruent simple nodes, while the divert phase diverts all edges of their origins to
a single node, rendering all other nodes dead. Both phases of Algorithm~\ref{alg:cne} perform a
simple top-down traversal, recursively processing subregions of structural nodes
annotating inputs, outputs, arguments, and results, as well as simple nodes as congruent. For
\gammanodes, the algorithm marks only computations within a \emph{single} region as congruent and
performs no analysis between regions. In the case of \thetanodes, computations are only congruent
when they are congruent before and after the loop execution, \ie, the inputs and results of two
loop variables must be congruent.
Figure~\ref{fig:dcne1} shows the RVSDG for the code in Figure~\ref{fig:dcne-code}, and Figure
\ref{fig:dcne2} the RVSDG after CNE. Two of the four multiplications take the same inputs and
therefore are congruent to each other, resulting in the redirection of their result edges.

\begin{informalg}{Common Node Elimination}{alg:cne}
\begin{enumerate}
	\vspace*{-0.7\baselineskip}

	\item \bpoint{Mark}: Process all nodes in topological order as follows:
		\begin{itemize}
			\item \bpoint{Simple nodes}: Denote this node as \n. Mark \n as congruent to all nodes
				\np which represent the same operation and where $\card{I_n} = \card{I_{n^{\prime}}} \land
				\iput{k}{n} \cong \iput{k}{n^{\prime}}$ for all $k = [0..\card{I_n}]$. Mark all outputs
				$\oput{k}{n} \cong \oput{k}{n^{\prime}}$ for all $k = [0..\card{O_n}]$.

			\item \bpoint{\gammanode}: For all entry variables $ev_1, ev_2 \in EV$ where
				$\iput{}{ev_1} \cong \iput{}{ev_2}$, mark $\argument{k}{ev_1} \cong \argument{k}{ev_2}$ for
				all $k \in [0..\card{A_{ev_1}}]$. Recursively process the \gammaregions. For all exit
				variables $ex_1, ex_2 \in EX$ where $\result{k}{ex_1} \cong \result{k}{ex_2}$ for all
				$k \in [0..\card{R_{ex_1}}]$, mark $\oput{}{ex_1} \cong \oput{}{ex_2}$.

			\item \bpoint{\thetanode}: For all loop variables $lv_1, lv_2 \in LV$ where
				$\iput{}{lv_1} \cong \iput{}{lv_2} \land \result{}{lv_1} \cong \result{}{lv_2}$, mark
				$\argument{}{lv_1} \cong \argument{}{lv_2}$ and $\oput{}{lv_1} \cong \oput{}{lv_2}$.
				Recursively process the \thetaregion.

			\item \bpoint{\lambdanode}: For all context variables $cv_1, cv_2 \in CV$ where
				$\iput{}{cv_1} \cong \iput{}{cv_2}$, mark $\argument{}{cv_1} \cong \argument{}{cv_2}$.
				Recursively process the \lambdaregion.

			\item \bpoint{\phinode}: For all context variables $cv_1, cv_2 \in CV$ where
				$\iput{}{cv_1} \cong \iput{}{cv_2}$, mark $a_{cv_1} \cong a_{cv_2}$. Recursively process
				the \phiregion.

			\item \bpoint{\omeganode}: Recursively process the \omegaregion.
		\end{itemize}

	\item \bpoint{Divert}: Process all nodes in topological order as follows:
		\begin{itemize}
			\item \bpoint{Simple nodes}: Denote this node as \n. For all nodes \np which are congruent to
				\n, divert all outputs $\oput{k}{n^{\prime}}$ to \oput{k}{n} for all $k = [0..\card{O_n}]$.

			\item \bpoint{\gammanode}: For all entry variables $ev_1, ev_2 \in EV$ where
				$\iput{}{ev_1} \cong \iput{}{ev_2}$, divert all edges from $\argument{k}{ev_2}$ to
				$\argument{k}{ev_1}$ for all $k \in [0..\card{A_{ev_1}}]$. Recursively
				process the \gammaregions. For all exit variables $ex_1, ex_2 \in EX$ where
				$\result{k}{ex_1} \cong \result{k}{ex_2}$ for all $k \in [0..\card{R_{ex_1}}]$, divert all
				edges from $\oput{}{ex_2}$ to $\oput{}{ex_1}$.

			\item \bpoint{\thetanode}: For all induction variables $lv_1, lv_2 \in LV$ where
				$\argument{}{lv_1} \cong \argument{}{lv_2} \land \oput{}{lv_1} \cong \oput{}{lv_2}$, divert
				all edges from $\argument{}{lv_2}$ to $\argument{}{lv_1}$ and from $\oput{}{lv_2}$ to
				$\oput{}{lv_1}$. Recursively process the \thetaregion.

			\item \bpoint{\lambdanode}: For all context variables $cv_1, cv_2 \in CV$ where
				$\iput{}{cv_1} \cong \iput{}{cv_2}$, divert all edges from $\argument{}{cv_2}$ to
				$\argument{}{cv_1}$. Recursively process the \lambdaregion.

			\item \bpoint{\phinode}: For all context variables $cv_1, cv_2 \in CV$ where
				$\iput{}{cv_1} \cong \iput{}{cv_2}$, divert all edges from $\argument{}{cv_2}$ to
				$\argument{}{cv_1}$. Recursively process the \phiregion.

			\item \bpoint{\omeganode}: Recursively process the \omegaregion.
		\end{itemize}

	\vspace*{-\baselineskip}
\end{enumerate}
\end{informalg}

For simple nodes, the algorithm marks all nodes within a
region that are congruent to a node \n. In order to avoid costly traversals of all nodes for every
node \n, the mark phase takes the candidates from the users of the origin of \n's first
input. If there is another input from a simple node \np with the same operation and number of
inputs among them, the other inputs from both nodes can be compared for congruence. Moreover, a
region must store constant nodes, \ie nodes without inputs, separately from other nodes so that the
candidate nodes for constants are available. For commutative simple nodes, the inputs should be
sorted before their comparison.

The presented algorithm only detects congruent simple nodes within a region. For \gammanodes,
congruence can also exist between nodes of different \gammaregions, and extending the
algorithm would eliminate these redundancies. Another extension would
be to detect congruent structural nodes, to implement conditional fusion
\cite{Rugina00} and loop fusion~\cite{manjikian97}. In
the case of \gammanodes, it is sufficient to ensure that two nodes have congruent
predicates, while \thetanodes require congruence detection between
different \thetaregions to ensure that their predicates are the same.

